\begin{theorem}[The compact nodal plumbing]\hypertarget{tag-02ZI}{}\label{push:src:A.1.2}\label{push:A12}
\StatementRecord{push:src:A.1.2}{proved}{02ZI}{\SeriesRef{push:A11}; \cite{EGFS25}}
\hypertarget{tag-01R2}{}\label{fam:src:E.1.2}\label{fam:st:E.1.2}
\StatementRecord{fam:src:E.1.2}{proved}{01R2}{\SeriesRef{fam:src:E.1.1}; \cite{Engel2026}}

There is a reduced complex space $\mathcal X^\flat$ and a proper flat
holomorphic map
\[
 (f,\varpi):\mathcal X^\flat\longrightarrow B\times\Delta_u
\]
whose central threefold is $Q$. For $u\ne0$ its fibre is the
contraction of Engel's chosen $X(u)$ at one rational curve of the
Mumford fibre. This curve is the double curve corresponding to the
diagonal of the square in the chosen $A_2$ triangulation; in the
coordinates of the marked cusp fan its direction is $e_1-e_2$. Its normal bundle
in $X(u)$ is $\mathcal O_{\mathbb P^1}(-1)^{\oplus2}$. Its image is
the unique ordinary double point of $\mathcal X^\flat_u$.

Over $S^\circ=S\setminus S_\infty$ the total plumbing space is smooth. Its restriction over every relatively compact disc $U\Subset B\setminus\{\infty\}$ is proper and flat over $U\times\Delta_\epsilon$ for sufficiently small $\epsilon>0$. For $u\ne0$ it is the untwisted quotient by $u\phi$ covering $T^{-1}$. This includes discs containing $0$ or $1$, before the logarithmic transformations.
\end{theorem}


\begin{proof}
\begin{enumerate}[label=\arabic*.,leftmargin=*]
\item Here is a construction on the whole compact surface $S$. Choose a
Hermitian metric on $L$, use the dual metric on $L^{-1}$, and choose
$0<r<1/4$. For sufficiently small $\epsilon>0$ we have
$|u|\|\phi_x\|<r^2$ for all $x\in S$ and $|u|<\epsilon$.
Take the core
\[
 C=\{(x,a,u):a\in L_x,\ r<\|a\|<r^{-1},\ |u|<\epsilon\}
\]
and the node space
\begin{equation}
 N=\{(x,a,b,u):a\in L_x,\ b\in L^{-1}_{Tx},\
 a\otimes b=u\phi_x,\ \|a\|,\|b\|<2r,\ |u|<\epsilon\}.
\label{pub:I:eq:E1}
\end{equation}
On $r<\|a\|<2r$ identify a point of $N$ with $(x,a,u)\in C$.
On $r<\|b\|<2r$ identify it with $(Tx,b^{-1},u)\in C$.
Both identifications are biholomorphisms between open complex spaces;
in the first collar $b=u\phi_x/a$, and in the second
$a=u\phi_xb^{-1}$. Their domains are disjoint because
$\|a\|\|b\|<r^2$. Changes of frames multiply the equation by
units and give the stated bundle-valued equation, so no choice of
frame enters the gluing.

\item This gluing is Hausdorff. At either end of a collar, exactly one
of its two charts contains the limiting circle: the core omits
$\|a\|=r$, and the node chart omits $\|a\|=2r$; the infinity
collar has the identical property in $b$. Thus the graph of the
identification has no unpaired pair of limiting points. The two
collars in each chart are separated, so there are no additional
chains of identifications. These observations give disjoint
saturated neighbourhoods of any two inequivalent points. Properness
is equally concrete. Over a compact subset of $B\times\Delta_\epsilon$,
the closed core $3r/2\leq\|a\|\leq2/(3r)$ together with the
closed node subspace $\|a\|,\|b\|\leq7r/4$ is compact and maps
onto the glued space. To see surjectivity, a point outside the smaller
core lies on a node collar, and a node point outside the smaller
node subspace lies on the smaller core. Compactness of $S$ supplies
the uniform bound used above. This proves properness of the joint
map to $(t,u)$.

\item The total node rings are hypersurface rings over the regular
two-dimensional local rings of $S$. They are Cohen--Macaulay.
Every fibre of the joint map has pure dimension two: first fix
$t$, and then impose $xy=u\phi$ on the two line coordinates over
the one-dimensional fibre of $S$. Consequently $t-t_0,u-u_0$ is a
regular sequence in every total local ring. The regular-sequence
flatness criterion over the smooth two-dimensional base proves
flatness. On $S^\circ$, the derivative of $ab-u\phi_x$ with respect to $u$ is $-\phi_x\ne0$, so the total node space is smooth. On the core the total space is an open subset of $L\times\Delta_u$
and is smooth; joint flatness follows from that of $S\to B$. This includes the nonreduced Kodaira fibres
at $0,1$, since the argument uses smoothness of $S$. At $u=0$ the node equation is $ab=0$,
and the two collar maps recover precisely the specified pinching
of the two sections of $R$.

\item For $u\ne0$ and $t\ne\infty$, the equation gives
$a_{Tx}=a_x/(u\phi_x)$. Hence the resulting quotient is exactly
$L^\times/\langle u\phi\rangle$, where the contracting generator
covers $T^{-1}$. Every orbit meets a core or a collar, by the
uniform contraction estimate; the displayed comparison identifies
exactly its orbit relation. It is therefore a biholomorphism.

\item At the nodal point of $S_\infty$, write $q=zv$. Absorbing a unit
in one line coordinate gives $xy=uzv$. At a smooth point of
$S_\infty$ the equation is $xy=uq$. For $u\ne0$ the latter
threefold is smooth; the former has one nondegenerate quadratic
singularity at the origin. There are no other singularities of a
nonzero $u$ fibre.

\item In a Tate covering of a neighbourhood of the $I_1$
fibre use the torus coordinates $(Z,W)$ of
\citet[Proposition~2.7, p.~6]{Engel2026}.
The two periods act as
\[
 (Z,W,q,u)\longmapsto(c_1qZ,c_2W,q,u),\qquad
 (Z,W,q,u)\longmapsto(c_3Z,c'_4uqW,q,u),
\]
with holomorphic units $c_i$. Their monomial parts are $G_1,G_2$.
The chart indexed by $(k,l)$ has
$z=Zq^{-k}$, $v=q^{k+1}/Z$, $x=W(qu)^{-l}$,
$y=qu/x$, with the unit changes given by the line
frames. It is the chart of Proposition~\ref{push:A11}.
The Tate cover unfolds the node of $S_\infty$ into its infinite
chain, and unfolding the two plumbing collars gives the additional
integer index $l$. The foregoing collar identifications are exactly
the two deck identifications. Thus, on an invariant neighbourhood
over a sufficiently small bidisc, the fan space maps by a
holomorphic covering to the plumbing, with deck group $\mathbb Z^2$.
Each core or node coordinate neighbourhood has disjoint lifts
indexed by the two unfolded integers; this proves freeness and
proper discontinuity, including the $\rho_l$ orbits. The first
period stabilizes such a ray but on its orbit translates the Tate
coordinate by $c_1q$, with $q\ne0$ and $|c_1q|<1$; it does not
fix a point. Properness of its quotient was proved by the compact
core and node representatives. This proves the analytic comparison of the fan quotient with the compact plumbing.

\item Subdividing the square at $h_u=0$ by the chosen diagonal gives
the small resolution of $xy=uzv$ for $u\ne0$. Put $Y=y/u$, so the equation is $xY=zv$.
The resolution associated with the diagonal $[e_1,e_2]$ has charts
\[
 (v,Y,a)\longmapsto(av,uY,aY,v),\qquad
 (x,z,b)\longmapsto(x,ubz,z,bx)
\]
in the coordinates $(x,y,z,v)$. On their overlap,
$b=a^{-1}$ and $(x,z)=a(v,Y)$. The character $a=x/v$ has
valuation $y_1+y_2-h_q$, so these are exactly the two triangles
separated by $[e_1,e_2]$. The exceptional curve is the zero
section in the two normal coordinates. Their transition is
multiplication by $a$, giving
$\mathcal O_{\mathbb P^1}(-1)\oplus\mathcal O_{\mathbb P^1}(-1)$. There is one
diagonal edge modulo the periods, so there is exactly one such
curve. Removing that edge contracts the paired opposite
$(-1)$-curves in the normalization $dP_6$ to the square model
$\mathbb P^1\times\mathbb P^1$. This identifies the contracted
curve in the selected toric model. The fixed double curve of
Proposition~\ref{mapsI:src:A.1.2} has direction $e_2$ in the same
marked lattice; the two primitive directions satisfy
$\det(e_1-e_2,e_2)=1$.
\end{enumerate}
\end{proof}
